\section{بخش چهارم: مقایسه‌ی کیفی و بررسی زمان اجرا}

\begin{stepbox}
\field{مقایسه‌ی کیفی ۳ روش بخش‌بندی}{
بر اساس پیاده‌سازی‌های انجام‌شده بر روی تصاویر حیوانات در طبیعت، مقایسه‌ی کیفی الگوریتم‌ها به شرح زیر است:
\begin{itemize}
    \item \textbf{آستانه‌گذاری (\lr{Thresholding}):} ضعیف‌ترین عملکرد را دارد. به دلیل هم‌پوشانی بالای مقادیر پیکسلی سوژه و پس‌زمینه (نور، سایه و بافت مشابه شن‌ها یا ابرها)، نمی‌توان با یک آستانه‌ی سراسری، حیوان را به درستی جدا کرد.
    \item \textbf{\lr{Mean Shift}:} بهترین تعادل بین کیفیت و عملکرد را ارائه می‌دهد. بافت‌های اضافی پس‌زمینه را به خوبی یکدست کرده و لبه‌های سوژه را با دقت مناسبی حفظ می‌کند (به ویژه با ابعاد پنجره‌ی مکانی بهینه که سایه‌های اضافی را حذف می‌کند).
    \item \textbf{\lr{Normalized Cut}:} از نظر تئوری مفهومی قدرتمند است، اما در عمل به شدت به تنظیمات اولیه (مانند تعداد سوپرپیکسل‌ها) وابسته بوده و ریسک خطا در ادغام نواحی تصویر در آن بالاست.
\end{itemize}
}
\end{stepbox}

\vspace{0.5cm}
\newpage
\begin{stepbox}
\field{مکتوب‌سازی زمان‌های اجرا}{
زمان اجرای الگوریتم \lr{Normalized Cut} رابطه‌ی مستقیمی با ابعاد و رزولوشن تصویر دارد. در جدول زیر، زمان اجرای این الگوریتم برای سه تصویر با پارامترهای مختلف ثبت شده است. 
}
\end{stepbox}

% قرار دادن جدول در مرکز و جلوگیری از شکستن آن
\begin{table}[h!]
    \centering
    \caption{زمان اجرای الگوریتم \lr{Normalized Cut} به ازای مقادیر مختلف پارامترها (بر حسب ثانیه)}
    \label{tab:execution_times}
    \renewcommand{\arraystretch}{1.4}
    \begin{tabular}{p{4.5cm} c c c}
        \toprule
        \textbf{تصویر (حجم تقریبی فایل)} & \textbf{$c=10, n=100$} & \textbf{$c=30, n=200$} & \textbf{$c=50, n=400$} \\
        \midrule
        اسب ($2.2$ مگابایت) & $71.29$ & $73.40$ & $77.07$ \\
        ماهی ($800$ کیلوبایت) & $7.30$ & $8.19$ & $12.05$ \\
        پلنگ سیاه ($600$ کیلوبایت) & $10.47$ & $10.09$ & $11.59$ \\
        \bottomrule
    \end{tabular}
\end{table}

\vspace{0.5cm}

% برای جلوگیری از شکستن باکس روی مرز صفحه، آن را نشکن می‌کنیم
\begin{notebox}
همان‌طور که در جدول \ref{tab:execution_times} مشاهده می‌شود، به دلیل بالا بودن رزولوشن تصویر اسب نسبت به دو تصویر دیگر، مرتبه‌ی زمانی اجرای الگوریتم مبتنی بر گراف به صورت تصاعدی افزایش یافته و از میانگین $10$ ثانیه برای تصاویر کوچکتر، به بیش از $70$ ثانیه رسیده است. این در حالی است که سرعت اجرای روش آستانه‌گذاری در حد کسری از ثانیه و الگوریتم \lr{Mean Shift} در حدود چند ثانیه ارزیابی شد.
\end{notebox}